\documentclass[12pt,a4paper]{article}

\usepackage[margin=2.5cm]{geometry}
\usepackage{amsmath,amssymb}
\usepackage{graphicx}
\usepackage{booktabs}
\usepackage{hyperref}
\usepackage{xcolor}
\usepackage{listings}
\usepackage{caption}
\usepackage{float}
\usepackage{parskip}
\usepackage{titlesec}
\usepackage{enumitem}
\usepackage{multirow}
\usepackage{array}
\usepackage{fancyhdr}
\usepackage{microtype}
\usepackage{setspace}

\onehalfspacing

\hypersetup{
  colorlinks=true,
  linkcolor=blue!70!black,
  citecolor=blue!70!black,
  urlcolor=blue!70!black,
}

\lstset{
  basicstyle=\ttfamily\footnotesize,
  breaklines=true,
  frame=single,
  backgroundcolor=\color{gray!8},
}

\setlength{\headheight}{15pt}
\pagestyle{fancy}
\fancyhf{}
\rhead{Adaptive Quiz ML Pipeline}
\lhead{Yahia Kerroum -- ENSIA ML Spring 2026}
\cfoot{\thepage}

\titleformat{\section}{\large\bfseries}{}{0em}{}[\titlerule]
\titleformat{\subsection}{\normalsize\bfseries}{}{0em}{}

\title{
  \vspace{-1cm}
  {\Large\bfseries Adaptive Quiz ML Pipeline} \\[0.5em]
  {\large An End-to-End Adaptive Assessment System\\
  Using Active Learning and Gradient Boosting} \\[1em]
  {\normalsize Machine Learning Course -- ENSIA, Spring 2025--2026}
}
\author{Yahia Kerroum}
\date{May 2026}

\begin{document}

\maketitle
\thispagestyle{fancy}

\begin{abstract}
We design, implement, and evaluate an end-to-end adaptive quiz system that classifies
student knowledge levels (\emph{beginner}, \emph{intermediate}, \emph{advanced}) and
selects questions dynamically to minimise the number of questions required while
maximising classification accuracy.
The pipeline proceeds through nine stages: question bank curation (8,561 questions across
10 modules), synthetic student generation under a realistic $\mathcal{N}(0.55,\,0.18)$
population model, 25-feature engineering, multi-module Logistic Regression training
(85.7\% grouped CV), adaptive strategy comparison (QBC achieves a 40\% question
reduction), and four real-data validation experiments.
A purpose-built single-module model is then developed iteratively: starting from a
Logistic Regression baseline (\(\approx\)71\% CV), we upgrade to
\textbf{HistGradientBoosting} with 21 discriminative features (v3: 81.2\% CV) and
finally incorporate population-realistic class proportions and support for 8-question
sessions (v4: 77.0\% CV -- a more honest metric on a real student distribution).
The deployed backend integrates length-conditional confidence thresholds that allow
early stopping as soon as question 8 for high-confidence students.
\end{abstract}

\tableofcontents
\newpage

% ============================================================================
\section{Presentation Roadmap}
\label{sec:roadmap}
% ============================================================================

This report is structured to mirror the 15-minute presentation format required for
the final evaluation. Each major section below maps to one of the five presentation
pillars:

\begin{table}[H]
\centering
\caption{Report sections and presentation mapping}
\begin{tabular}{@{}p{4.5cm}p{4.5cm}p{5cm}@{}}
\toprule
Presentation point & Report section & Key claim \\
\midrule
Problem Statement / Context &
\S\ref{sec:intro} &
Classify student knowledge level using only answer sequences; minimise questions needed. \\[4pt]
Data Description \& Preprocessing &
\S\ref{sec:data}, \S\ref{sec:population} &
8,561 MCQs from Sanfoundry; hybrid difficulty labelling; synthetic $\mathcal{N}(0.55,0.18)$ population. \\[4pt]
Solution Design &
\S\ref{sec:features}, \S\ref{sec:multimodule}, \S\ref{sec:singlemodule} &
25/21-feature vectors; HistGBM with variance features; population-realistic training. \\[4pt]
Experimentation \& Results Comparison &
\S\ref{sec:adaptive}, \S\ref{sec:realdata}, \S\ref{sec:singlemodule} &
QBC 40\% reduction; 77\%--96\% CV across datasets; model evolution LogReg $\to$ HistGBM. \\[4pt]
Conclusion &
\S\ref{sec:conclusion}, \S\ref{sec:discoveries} &
Working deployed platform; identified structural limits of single-module classification. \\
\bottomrule
\end{tabular}
\end{table}

% ============================================================================
\section{Introduction}
\label{sec:intro}
% ============================================================================

Adaptive testing, or Computer-Adaptive Testing (CAT), adjusts the difficulty and
selection of questions in real time based on a student's emerging performance profile.
Unlike fixed-form quizzes where every student answers the same questions in the same
order, adaptive systems can reach the same classification confidence with significantly
fewer questions -- reducing assessment fatigue and improving the signal-to-noise ratio
for borderline students.

This project builds a fully functional adaptive quiz platform, integrating a machine
learning classification backend with a Next.js frontend and a FastAPI/PostgreSQL
backend. The ML component is responsible for three tasks:
\begin{enumerate}[noitemsep]
  \item \textbf{Level prediction}: classify a student as beginner, intermediate, or
        advanced given their answer history.
  \item \textbf{Question selection}: given remaining candidates, choose the question
        whose difficulty tier best targets the current uncertainty.
  \item \textbf{Stop criterion}: decide when enough evidence has accumulated to commit
        to a classification, terminating the session early.
\end{enumerate}

The report is structured as follows.
Section~\ref{sec:data} describes the question bank and difficulty labelling.
Section~\ref{sec:population} covers synthetic student generation.
Section~\ref{sec:features} details the feature engineering pipeline.
Section~\ref{sec:multimodule} presents the multi-module model.
Section~\ref{sec:adaptive} evaluates adaptive strategies.
Section~\ref{sec:realdata} summarises real-data validation.
Section~\ref{sec:singlemodule} traces the evolution of the single-module model.
Section~\ref{sec:thresholds} derives the adaptive confidence thresholds.
Section~\ref{sec:platform} describes the platform integration.
Section~\ref{sec:limitations} discusses limitations and Section~\ref{sec:conclusion}
concludes.

% ============================================================================
\section{Dataset and Question Bank Curation}
\label{sec:data}
% ============================================================================

\subsection{Raw data source}

The question bank is scraped from Sanfoundry (\url{https://www.sanfoundry.com}), a
large repository of multiple-choice questions for engineering topics.
The raw bank contains \textbf{8,561 questions} spanning 10 computer science modules:
Computer Fundamentals, C++ Programming, Cyber Security, Data Structures,
Operating Systems, Software Engineering, and four additional subjects.

\subsection{Three-stage filtering}

Questions pass three consecutive filters before entering the adaptive pipeline:

\begin{table}[H]
\centering
\caption{Question bank filtering pipeline}
\begin{tabular}{@{}llrr@{}}
\toprule
Stage & Filter & Removed & Remaining \\
\midrule
1 & Module selection (6 core subjects) & 2,200 & 6,361 \\
2 & Quality flags (duplicate, malformed) & 1,006 & 5,355 \\
3 & Semantic deduplication (cosine > 0.85) & 997  & \textbf{4,358} \\
\bottomrule
\end{tabular}
\end{table}

Semantic deduplication uses TF-IDF embeddings with cosine similarity; any question
pair with similarity above 0.85 has the shorter question removed, retaining the more
complete formulation.

\subsection{Difficulty labelling}

Difficulty labels are assigned with a hybrid strategy:
\begin{enumerate}[noitemsep]
  \item \textbf{Gemini-assisted labelling} (primary): a structured prompt asks the
        model to classify each question as easy/medium/hard with a confidence score.
        Questions with confidence $\geq 0.80$ receive their label directly.
  \item \textbf{Manual ground-truth} (secondary): 150 questions per module are
        hand-labelled by the project author; these serve as the reference set.
  \item \textbf{Rubric scoring} (fallback): questions failing both above criteria
        receive a rule-based score based on question length, number of distractors,
        and lexical complexity.
\end{enumerate}

The final distribution across the 6 retained modules is approximately
35\% easy, 42\% medium, 23\% hard -- closely matching the synthetic session
difficulty distribution used during training.

% ============================================================================
\section{Synthetic Student Population}
\label{sec:population}
% ============================================================================

\subsection{Population model}

Real student data is unavailable at training time; the model is therefore trained
entirely on \emph{synthetic} sessions. The population model is parameterised by a
Normal distribution $\mathcal{N}(\mu,\sigma)$ over overall accuracy:

\[
  \text{overall\_acc} \sim \mathcal{N}(0.55,\;0.18)
\]

This places the average student at 55\% accuracy with a standard deviation of 18
percentage points. Integrating under the Normal curve with class boundaries at
$\text{acc} = 0.45$ (beginner/intermediate) and $\text{acc} = 0.70$
(intermediate/advanced) yields the following \textbf{realistic class proportions}:

\begin{table}[H]
\centering
\caption{Class proportions from $\mathcal{N}(0.55,\,0.18)$}
\begin{tabular}{@{}lcc@{}}
\toprule
Class & Accuracy range & Population fraction \\
\midrule
Beginner     & $< 0.45$          & $\approx 29\%$ \\
Intermediate & $[0.45,\;0.70)$   & $\approx 51\%$ \\
Advanced     & $\geq 0.70$       & $\approx 20\%$ \\
\bottomrule
\end{tabular}
\end{table}

\subsection{Per-level generative model}

Each student's session is generated from a \emph{per-difficulty accuracy profile}
sampled around level-specific centres:

\begin{table}[H]
\centering
\caption{Level centres (accuracy by difficulty)}
\begin{tabular}{@{}lccc@{}}
\toprule
Level & Easy & Medium & Hard \\
\midrule
Advanced     & 0.94 & 0.82 & 0.68 \\
Intermediate & 0.74 & 0.54 & 0.30 \\
Beginner     & 0.50 & 0.27 & 0.13 \\
\bottomrule
\end{tabular}
\end{table}

A per-session profile is sampled: $p_d \sim \mathcal{N}(\mu_d, 0.055)$, clipped to
$[0.04, 0.98]$, with the monotonicity constraint $p_\text{easy} \geq p_\text{medium}
\geq p_\text{hard}$ enforced. Response times follow independent Normal distributions
per level and difficulty, calibrated so that beginners guess quickly on hard questions
while advanced students slow down on them.

Boundary students (15\% of each level) have their accuracy centres shifted toward
the adjacent class to simulate realistic overlap in the classifier's decision regions.

% ============================================================================
\section{Feature Engineering}
\label{sec:features}
% ============================================================================

Each student session is converted into a fixed-length feature vector.
The multi-module model uses 25 features; the final single-module model uses 21.
All indices 0--14 are identical across both models, ensuring backward compatibility
in the backend service.

\subsection{Feature groups}

\begin{table}[H]
\centering
\caption{21-feature vector for the single-module model}
\begin{tabular}{@{}clp{7cm}@{}}
\toprule
Index & Feature & Description \\
\midrule
0     & \texttt{overall\_acc}       & Fraction of correct answers \\
1--3  & \texttt{acc\_easy/med/hard} & Per-difficulty accuracy \\
4--6  & \texttt{has\_easy/med/hard} & Binary: difficulty tier seen \\
7     & \texttt{last3\_acc}         & Accuracy on last 3 questions (recency) \\
8     & \texttt{error\_streak}      & Current consecutive error count \\
9     & \texttt{correct\_streak}    & Current consecutive correct count \\
10    & \texttt{acc\_trend}         & \texttt{last3\_acc} $-$ \texttt{first3\_acc} \\
11    & \texttt{fast\_correct\_rate}& Fraction fast \& correct \\
12    & \texttt{fast\_wrong\_rate}  & Fraction fast \& wrong (guessing signal) \\
13    & \texttt{slow\_correct\_rate}& Fraction slow \& correct \\
14    & \texttt{time\_acc\_corr}    & Pearson $r$(speed, correctness), z-scored \\
15    & \texttt{weighted\_acc}      & $\frac{\sum w_d \cdot n_d \cdot \text{acc}_d}{\sum w_d n_d}$, $w=\{1,2,3\}$ \\
16    & \texttt{n\_norm}            & Session length / 20 \\
17    & \texttt{hard\_variance}     & Variance of correctness on hard questions \\
18    & \texttt{medium\_variance}   & Variance of correctness on medium questions \\
19    & \texttt{medium\_hard\_gap}  & $\text{acc\_medium} - \text{acc\_hard}$ \\
20    & \texttt{module\_id\_enc}    & Module identity, categorical (0--5) \\
\bottomrule
\end{tabular}
\end{table}

\subsection{Self-normalised time features}

Time-accuracy interaction features (indices 11--14) are \emph{self-normalised}: the
threshold between ``fast'' and ``slow'' is the student's own session mean, not an
absolute value. This makes features invariant to total quiz duration, device speed,
and inter-student differences in typing or reading pace.

\subsection{Variance and gap features}

Features 17--19 were the key addition that resolved the intermediate-class recall
collapse. Intermediate students are characterised by \emph{inconsistent} performance
on hard questions -- sometimes correct, sometimes wrong -- giving \texttt{hard\_variance}
values near 0.25 (maximum for a Bernoulli variable). Beginners have uniformly near-zero
hard accuracy (variance $\approx$ 0) and advanced students are uniformly near-1
(variance $\approx$ 0). The medium-hard gap similarly discriminates: advanced students
maintain a small gap ($\approx 0.14$), while intermediates show a large gap ($\approx
0.24$) and beginners show a very small gap ($\approx 0.14$ from low base rates).

% ============================================================================
\section{Multi-Module Model}
\label{sec:multimodule}
% ============================================================================

\subsection{Architecture}

The multi-module model classifies sessions in which the student answers 30 questions
spanning all 6 modules. Four classifiers are compared using 5-fold
\texttt{StratifiedGroupKFold} (grouped by student to prevent data leakage):

\begin{table}[H]
\centering
\caption{Classifier comparison on 25-feature multi-module data}
\begin{tabular}{@{}lcc@{}}
\toprule
Classifier & Grouped 5-fold CV & Notes \\
\midrule
Gaussian Naive Bayes   & 74.2\% & Fast but assumes feature independence \\
Decision Tree (depth 6)& 79.1\% & Interpretable; prone to overfitting \\
Random Forest          & 83.4\% & Ensemble, slow at inference \\
\textbf{Logistic Regression} & \textbf{85.7\%} & Best balance of accuracy and speed \\
\bottomrule
\end{tabular}
\end{table}

Logistic Regression is selected for the multi-module setting: it achieves the highest
CV accuracy and provides well-calibrated probability outputs necessary for the
confidence-based stop criterion.

\subsection{Feature ablation}

Removing the two time features (\texttt{avg\_time}, \texttt{time\_trend}) reduces CV
accuracy from 85.7\% to 83.1\% ($\Delta = -2.6$ pp), confirming that response-time
information carries independent predictive signal beyond raw accuracy.

\subsection{Calibration}

Reliability diagrams show that the raw Logistic Regression is slightly overconfident
in the 60--80\% range. Platt scaling (sigmoid calibration via cross-validated
\texttt{CalibratedClassifierCV}) corrects this, producing predicted probabilities
that track empirical frequencies within $\pm 3\%$ across all bins.

% ============================================================================
\section{Adaptive Strategies}
\label{sec:adaptive}
% ============================================================================

Four adaptive question-selection strategies are evaluated in simulation:

\begin{enumerate}[noitemsep]
  \item \textbf{Random}: uniform random selection from remaining questions.
  \item \textbf{Uncertainty sampling}: select the question whose difficulty tier
        maximises classifier uncertainty (smallest max-probability class).
  \item \textbf{Entropy}: select the question that would maximally increase
        predictive entropy if answered correctly vs incorrectly.
  \item \textbf{Query-by-Committee (QBC)}: a committee of three classifiers
        votes; select the question with maximum vote disagreement.
\end{enumerate}

\subsection{Results}

\begin{table}[H]
\centering
\caption{Adaptive strategy comparison (500 simulated students)}
\begin{tabular}{@{}lccc@{}}
\toprule
Strategy & AUQC & Questions for 70\% accuracy & Final accuracy \\
\midrule
Random       & 0.712 & 10 & 84.1\% \\
Uncertainty  & 0.718 & 8  & 84.9\% \\
Entropy      & 0.721 & 8  & 85.2\% \\
\textbf{QBC} & \textbf{0.730} & \textbf{6} & \textbf{85.7\%} \\
\bottomrule
\end{tabular}
\end{table}

QBC achieves a \textbf{40\% question reduction} (6 questions vs 10 for random) to
reach 70\% population accuracy. The Area Under the Questions-to-Classification
curve (AUQC) summarises this efficiency across all thresholds.

\subsection{Statistical significance}

A paired $t$-test comparing QBC and random accuracy-vs-questions trajectories
yields $p < 0.001$, confirming the improvement is not due to sampling variance.
Per-module analysis shows the largest gains in Operating Systems and Data Structures
modules, where questions are more uniformly distributed across difficulties.

% ============================================================================
\section{Real Data Validation}
\label{sec:realdata}
% ============================================================================

The model is validated on four real datasets of increasing domain distance:

\begin{table}[H]
\centering
\caption{Cross-curriculum generalisation results}
\begin{tabular}{@{}lccl@{}}
\toprule
Dataset & Students & CV Accuracy & Notes \\
\midrule
Data Mining quizzes (ENSIA) & 127 & 89.4\% $\pm$ 4.1\% & Same institution, 5 quizzes \\
OS/ITE quizzes (AI224)      & 626 & 71.3\% $\pm$ 6.2\% & 92\% class imbalance \\
Lectu/ENSIA                 & 412 & 96.6\% $\pm$ 3.2\% & Real adaptive sessions \\
EdNet KT1/KT2               & 5,000 & 93.1\% $\pm$ 1.8\% & Large-scale, different domain \\
\bottomrule
\end{tabular}
\end{table}

The OS/ITE result (71.3\%) is lower due to severe class imbalance -- 578 of 626 students
are labelled intermediate, leaving only 48 advanced students for the rare class. The
high accuracy on Lectu/ENSIA and EdNet confirms that the model's features are
domain-agnostic: accuracy profiles and time-accuracy interactions transfer well
to unseen curricula.

\subsection{Online inflation}

Data Mining scores show a systematic 8--12 pp upward bias compared to synthetic
predictions. This is consistent with online quiz taking: students answer with access
to course notes, creating an inflated accuracy signal. The model correctly detects
these students as intermediate-to-advanced, which is accurate for their
\emph{assisted} performance even if their unassisted level is lower.

% ============================================================================
\section{Single-Module Model Evolution}
\label{sec:singlemodule}
% ============================================================================

The multi-module model requires 30 questions spanning 6 modules.
The live platform serves single-subject quizzes (10--20 questions from one module).
A dedicated single-module model is therefore developed.

\subsection{v1: Logistic Regression baseline}

The first attempt extends the multi-module feature set with 12 module-pair features
(one accuracy and one binary flag per module). The 26-feature LogReg reaches only
\textbf{$\approx$71\% CV}, with intermediate class recall collapsing to \textbf{35\%}.

Diagnosis: LogReg cannot learn the non-linear decision boundary between intermediate
and beginner/advanced. Intermediate students are characterised by \emph{inconsistency}
on hard questions -- a variance-based pattern invisible to a linear model.

\subsection{v2: Calibrated LogReg with grid search}

GridSearchCV over regularisation strength ($C \in \{0.5, 1, 2, 5, 10\}$) and solver
finds the optimum at $C=5$, \texttt{lbfgs}. Platt calibration is applied.
CV improves to $\approx$73\%, confirming that 71\% is a ceiling for linear models
on this feature set, not a tuning deficiency.

\subsection{v3: HistGradientBoosting, 21 features}
\label{sec:v3}

Three structural changes break through the linear ceiling:

\begin{enumerate}[noitemsep]
  \item \textbf{Feature reduction}: 12 module-pair features replaced by a single
        categorical \texttt{module\_id\_enc} (integer 0--5). HistGBM creates
        dedicated split bins for categorical features, recovering the module signal
        without the 12-feature redundancy.
  \item \textbf{Variance features}: \texttt{hard\_variance}, \texttt{medium\_variance},
        \texttt{medium\_hard\_gap} added (features 17--19). These are the primary
        discriminators for the intermediate class.
  \item \textbf{Classifier}: \texttt{HistGradientBoostingClassifier} with
        \texttt{categorical\_features=[20]}, trained on 12--20 question sessions
        (1500 per level/module, 15\% boundary students).
\end{enumerate}

\begin{table}[H]
\centering
\caption{HistGBM v3 hyperparameters}
\begin{tabular}{@{}ll@{}}
\toprule
Parameter & Value \\
\midrule
\texttt{max\_iter}          & 800 \\
\texttt{max\_depth}         & 7 \\
\texttt{learning\_rate}     & 0.03 \\
\texttt{min\_samples\_leaf} & 20 \\
\texttt{l2\_regularization} & 0.4 \\
\texttt{class\_weight}      & balanced \\
\bottomrule
\end{tabular}
\end{table}

\textbf{Result: 81.2\% CV $\pm$ 0.2\%}, intermediate recall rises from 35\% to
\textbf{86\%}. The low fold variance ($\pm$0.2\%) confirms structural improvement.

\subsection{v4: Realistic distribution, 8-question sessions}
\label{sec:v4}

Two problems remain after v3:
\begin{enumerate}[noitemsep]
  \item Sessions 12--20 questions forces \texttt{MIN\_QUESTIONS=12}. Most platform
        quizzes contain $\leq$15 questions, so adaptive stopping rarely activates
        before question 12.
  \item Equal class sampling (1500/level) is unrealistic and inflates the CV metric.
        The $\mathcal{N}(0.55, 0.18)$ population model predicts 51\% intermediate --
        the hardest class to classify.
\end{enumerate}

\textbf{Changes in v4:}
\begin{itemize}[noitemsep]
  \item Sessions 8--20 questions (40\% short 8--11, 60\% long 12--20).
  \item Class proportions: beginner 29\%, intermediate 51\%, advanced 20\%.
  \item Increased regularisation (\texttt{l2\_regularization}: 0.4 $\to$ 0.6;
        \texttt{min\_samples\_leaf}: 20 $\to$ 30) to handle noisier short sessions.
  \item More iterations (800 $\to$ 1000), shallower depth (7 $\to$ 6).
\end{itemize}

\begin{table}[H]
\centering
\caption{v4 accuracy by session length}
\begin{tabular}{@{}lcc@{}}
\toprule
Session length & Accuracy & Sessions \\
\midrule
n = 8--11  & 76.4\% & 10,794 \\
n = 12--15 & 78.8\% & 7,049 \\
n = 16--20 & 82.7\% & 9,151 \\
\midrule
\textbf{Overall CV} & \textbf{77.0\% $\pm$ 0.6\%} & 26,994 \\
\bottomrule
\end{tabular}
\end{table}

The 4.2 pp drop from v3 (81.2\% $\to$ 77.0\%) is expected and desirable: the CV is
now measured on the same imbalanced distribution as the real student population.
The 81.2\% was inflated by equal-class sampling; 77.0\% is the correct estimate of
real-world classification performance.

% ============================================================================
\section{Adaptive Confidence Thresholds}
\label{sec:thresholds}
% ============================================================================

\subsection{Motivation}

A flat confidence threshold (e.g., 80\%) treats a session of 8 questions the same as
one of 18. This is inappropriate: at $n=8$, each observation carries high binomial
variance. A student who correctly answers 6 of 8 hard questions could be advanced or
could have been lucky. A flat threshold would incorrectly terminate such sessions.

\subsection{Length-conditional design}

The threshold $\tau(n)$ is defined as:

\[
\tau(n) = \begin{cases}
  \infty    & n < 8  \quad\text{(never stop)} \\
  0.93      & 8 \leq n \leq 9 \\
  0.88      & 10 \leq n \leq 11 \\
  0.83      & 12 \leq n \leq 14 \\
  0.80      & n \geq 15
\end{cases}
\]

The thresholds are chosen to require near-certainty for early stops: a 93\%
confidence at $n=8$ corresponds to eliminating all but the most extreme student
profiles from early termination. The model's 88\% confidence on the canonical
advanced 8-question sanity check falls below 93\%, meaning that session correctly
continues to accumulate more evidence.

\subsection{Practical effect}

\begin{itemize}[noitemsep]
  \item \textbf{Clear beginners and advanced students} (flat low/high profiles)
        reach 93\%+ confidence by question 8--9 and stop early.
  \item \textbf{Intermediate and boundary students} rarely exceed 83\% at $n<12$,
        so they always answer at least 12 questions -- matching the regime where
        the model was most carefully validated.
  \item \textbf{Maximum length}: 20 questions; guaranteed termination.
\end{itemize}

% ============================================================================
\section{Platform Integration}
\label{sec:platform}
% ============================================================================

\subsection{Architecture}

The adaptive quiz platform is a full-stack web application:

\begin{itemize}[noitemsep]
  \item \textbf{Backend}: FastAPI (Python), async SQLAlchemy, PostgreSQL via asyncpg.
  \item \textbf{Frontend}: Next.js 14 (App Router), TypeScript, Tailwind CSS.
  \item \textbf{Auth}: Supabase JWT with email allowlist for admin access.
\end{itemize}

\subsection{ML service API}

The backend exposes four public functions in \texttt{backend/services/ml\_service.py}:

\begin{lstlisting}[language=Python]
compute_features(is_correct_list, difficulty_list,
                 time_ms_list, module_list) -> np.ndarray  # (1, 21)

predict_level(features) -> dict  # {level, confidence, probabilities}

select_next_question(features, candidate_numbers,
                     candidate_difficulties) -> int  # question number

should_stop(features, n_answered) -> bool
\end{lstlisting}

\texttt{should\_stop} integrates the length-conditional threshold:
\begin{lstlisting}[language=Python]
def should_stop(features, n_answered):
    if n_answered < MIN_QUESTIONS:  # MIN_QUESTIONS = 8
        return False
    if n_answered >= MAX_QUESTIONS:  # MAX_QUESTIONS = 20
        return True
    ...
    return prediction["confidence"] >= _confidence_threshold(n_answered)
\end{lstlisting}

\subsection{Model loading}

The model is loaded lazily on first call and cached globally. The bundle format is:
\begin{lstlisting}[language=Python]
bundle = {
    "model": HistGradientBoostingClassifier,
    "approach": "single_module_histgbm_v4",
    "n_features": 21,
    "feature_names": [...],
    "cv_accuracy": 0.770,
}
\end{lstlisting}

If the model file is missing, \texttt{predict\_level} falls back to a rule-based
heuristic using \texttt{acc\_hard} and \texttt{acc\_medium} thresholds, ensuring
the platform degrades gracefully.

\subsection{Session flow}

\begin{enumerate}[noitemsep]
  \item Student selects a quiz module and starts a session.
  \item The backend calls \texttt{select\_next\_question} with all remaining questions.
  \item After each answer, \texttt{compute\_features} is updated and \texttt{should\_stop}
        is evaluated.
  \item On termination, \texttt{predict\_level} returns the final classification;
        the session record is stored with level, confidence, and probabilities.
  \item The frontend displays a results page with the assigned level and score.
\end{enumerate}

% ============================================================================
\section{What We Discovered Along the Way}
\label{sec:discoveries}
% ============================================================================

This section documents non-obvious insights -- things that only became clear after
building and breaking the system. They are not visible from the final results alone.

\subsection{Why multi-module classification is easier than single-module}

The intuition is simple once stated: \emph{a strong student is strong in general}.
A student who scores 85\%+ on Operating Systems questions almost certainly also
performs well on Data Structures and Software Engineering. The multi-module feature
vector captures this by observing accuracy across six different knowledge domains in
a single session. The redundancy works in the classifier's favour: even if a student
has a bad run on one module, the signal from the other five corrects the estimate.

Single-module classification breaks this redundancy. A student who scores 75\% on
C++ might be genuinely intermediate, or might be an advanced student who never
touched C++ before, or a beginner who got lucky on easy questions. Without the
cross-module signal, there is no way to distinguish these cases from the single module
alone. This is why the multi-module model achieves 85.7\% CV while the single-module
model plateaus at 77\%: the problem is structurally harder, not just harder to tune.

The practical implication is that a platform offering cross-module quizzes would
get more accurate adaptive placement with fewer total questions than one that runs
strictly module-by-module -- a real system design trade-off.

\subsection{The structural problem with course-specific quiz data}

The OS/ITE dataset (626 students, AI224 course) delivers 71.3\% CV -- the worst
real-data result. The obvious explanation is data quality; the real explanation
is more fundamental.

The course delivers the same lectures to every student in a single semester cohort.
Students who enrol are, almost by construction, at a similar prior knowledge level:
they passed the prerequisite courses, they are in the same year of study, they are
preparing for the same exam. The result is that 578 of 626 students (92\%) end up
labelled intermediate -- not because the labelling is wrong, but because the
\emph{population is genuinely homogeneous}.

A classifier that sees 92\% of training examples in one class cannot learn to
distinguish the other two, regardless of how good the features are. This is not a
solvable data problem. It is a signal problem: the quiz was not designed to
discriminate among knowledge levels, it was designed to assess a specific curriculum.
Any adaptive system deployed in a similarly homogeneous context would face the same
ceiling, and 71\% with balanced class weights is close to the best possible result.

\subsection{The Data Mining quiz grade inflation problem}

The DM quizzes from ENSIA were administered online, via Google Forms, without
invigilation. Students answered with access to their course slides, textbooks, and
each other. The quiz scores are therefore \emph{assisted performance scores}, not
knowledge assessments.

This created a systematic upward shift of 8--12 percentage points compared to what
the synthetic model expects for each level. The model classifies the same student as
intermediate in the synthetic distribution but intermediate-to-advanced when given
their quiz scores -- not because it is wrong, but because it is accurately reflecting
the inflated input. The model cannot distinguish ``genuinely advanced'' from
``intermediate with open notes'': the feature vectors are identical.

The deeper issue is that this kind of data contamination is nearly impossible to
detect automatically. There is no feature in the answer history that reveals whether
a student consulted external resources. This is a fundamental limitation of any
online, unsupervised adaptive assessment -- and it is why supervised exam conditions
remain a hard requirement for high-stakes placement.

\subsection{Why we moved from raw response time to time-accuracy interaction}

The first version of the feature set included \texttt{avg\_time} and
\texttt{time\_trend} as raw millisecond values. These features have an obvious
psychometric motivation: students who think longer on hard questions are likely
more engaged than those who answer immediately. But raw times are confounded by
everything outside the student's knowledge: their device, their internet connection,
their reading speed, whether they misclicked, whether they stepped away from the
screen.

A student answering from a mobile phone over a slow connection will show systematically
higher times than the same student on a desktop, regardless of knowledge level.
The raw feature is a mixture of knowledge signal and device noise, and the model
cannot separate them.

The self-normalised approach resolves this by replacing ``is this response fast?''
with ``is this response fast \emph{for this student in this session}?'' The threshold
is the student's own session mean, so all device and connection effects cancel out.
What remains is the \emph{relative pattern}: did the student slow down on hard
questions (a sign of genuine engagement), or were their times uniform (suggesting
guessing)? The features \texttt{fast\_wrong\_rate} and \texttt{time\_acc\_corr}
capture this directly. The ablation study confirmed the improvement: removing the
time features drops CV by 2.6 pp.

\subsection{Why Logistic Regression hits a ceiling on intermediate classification}

When the intermediate class recall collapsed to 35\% with LogReg, the natural
response was to tune the regularisation parameter. GridSearchCV across five values
of $C$ moved accuracy from 71\% to 73\% -- a 2 pp improvement that exhausted the
entire tunable range. The ceiling was not a hyperparameter problem.

The root cause is geometric: intermediate students do not occupy a convex region in
feature space. An intermediate student is not someone who always gets medium accuracy
on every question type. They are someone whose hard-question performance is
\emph{unpredictable} -- sometimes correct, sometimes not -- while their easy-question
accuracy is reliably moderate. This inconsistency cannot be expressed as a linear
combination of features. It requires capturing variance, not just means.

HistGradientBoosting resolves this naturally because tree splits can express
``if hard\_variance is in this range AND acc\_medium is in that range, then
intermediate.'' Logistic Regression has no equivalent mechanism regardless of how
many features you add or how you regularise.

\subsection{The difficulty labelling gap}

The question bank contains 4,358 questions. Getting reliable difficulty labels for
all of them without a live student population to observe is genuinely hard. Our
approach -- Gemini-assisted labelling with manual ground truth for a sample -- works
well for extreme cases (very easy definition questions vs. very hard multi-step
reasoning problems) but struggles at the medium/hard boundary.

The practical consequence is that the ``medium'' difficulty bin is over-populated and
somewhat heterogeneous. Some questions labelled medium would be hard for a beginner and
easy for an advanced student, which is actually correct behaviour for an adaptive
system. But some medium questions are simply mislabelled, and the model has no way to
detect this: it treats them as accurate signals.

This is not a fixable problem without real student response data. The correct long-term
solution is to use the platform's own answer logs to estimate empirical difficulty
(fraction of correct responses per question) and retroactively update labels. This
is standard in industrial IRT (Item Response Theory) systems but requires a minimum
deployment period to accumulate sufficient data.

% ============================================================================
\section{Limitations and Honest Assessment}
\label{sec:limitations}
% ============================================================================

\begin{table}[H]
\centering
\caption{Known limitations}
\begin{tabular}{@{}p{3.5cm}p{3cm}p{7cm}@{}}
\toprule
Limitation & Severity & Impact and mitigation \\
\midrule
Synthetic training data only &
High &
The model has never observed a real student. The 77\% CV is measured on synthetic data;
real-world accuracy is unknown. Mitigation: collect labelled sessions as students use
the platform and retrain periodically. \\[4pt]
No cross-validation on real single-module sessions &
High &
The single-module model is validated via sanity checks, not held-out real data.
Mitigation: deploy with session logging to enable future validation. \\[4pt]
EdNet smoke test &
Medium &
EdNet contains different question types (knowledge tracing, binary responses only);
no ground-truth level labels are available. The 93.1\% CV uses proxy labels from
accuracy terciles, making it an upper-bound estimate. \\[4pt]
Intermediate class dominates &
Medium &
51\% of the population is intermediate under the $\mathcal{N}(0.55, 0.18)$ model.
Misclassifying intermediate students is the most common error (12--15\% error rate). \\[4pt]
Binomial noise at $n=8$ &
Low &
8 binary observations cannot support 95\%+ accuracy theoretically. The 93\%
threshold at $n=8$ partially mitigates this by only stopping extremely clear cases. \\[4pt]
Response time assumption &
Low &
Time features assume a minimum response of 500 ms. Platform response times often
exceed 5 s even for easy questions; absolute values are not used (self-normalised). \\
\bottomrule
\end{tabular}
\end{table}

% ============================================================================
\section{Conclusion}
\label{sec:conclusion}
% ============================================================================

This project delivers a complete adaptive quiz pipeline, from raw question scraping to
a deployed web application. The key technical contributions are:

\begin{enumerate}[noitemsep]
  \item A \textbf{21-feature vector} with self-normalised time-accuracy interaction
        and variance-based features that recover intermediate-class recall from 35\%
        to 86\%.
  \item A \textbf{HistGradientBoosting classifier} with a categorical module feature,
        outperforming Logistic Regression by 10 pp on the same data.
  \item \textbf{Population-realistic training data} derived from the
        $\mathcal{N}(0.55, 0.18)$ model, yielding a more honest 77\% CV metric on
        the actual student distribution.
  \item \textbf{Length-conditional confidence thresholds} that allow early stopping
        at $n=8$ for high-confidence students while ensuring ambiguous cases answer
        at least 12 questions.
  \item Cross-domain generalisation to four real datasets, with 93\%+ accuracy on
        Lectu/ENSIA and EdNet.
\end{enumerate}

The primary next step is replacing synthetic training data with real labelled sessions
as the platform accumulates usage. Even a small real-student dataset (200--300 sessions
per level) would allow fine-tuning the HistGBM and potentially push real-world accuracy
above the 77\% synthetic estimate.

% ============================================================================
\section*{Acknowledgements}
% ============================================================================

Question bank sourced from Sanfoundry (\url{https://www.sanfoundry.com}).
Real quiz data provided by ENSIA Data Mining and AI224 course instructors.
EdNet dataset provided by Riiid Inc.\ under the CC BY-NC 4.0 licence.

% ============================================================================
\begin{thebibliography}{9}

\bibitem{ednet}
  Choi, Y., et al.\ (2020).
  \emph{EdNet: A Large-Scale Hierarchical Dataset in Education}.
  AIED 2020.

\bibitem{histgbm}
  Biau, G., and Scornet, E.\ (2016).
  \emph{A Random Forest Guided Tour}.
  TEST, 25(2), 197--227.

\bibitem{irt}
  Lord, F.\ M.\ (1980).
  \emph{Applications of Item Response Theory to Practical Testing Problems}.
  Erlbaum.

\bibitem{cat}
  Wainer, H.\ (1990).
  \emph{Computerized Adaptive Testing: A Primer}.
  Erlbaum.

\bibitem{actlearn}
  Settles, B.\ (2012).
  \emph{Active Learning}.
  Morgan \& Claypool.

\bibitem{sklearn}
  Pedregosa, F., et al.\ (2011).
  \emph{Scikit-learn: Machine Learning in Python}.
  JMLR, 12, 2825--2830.

\bibitem{fastapi}
  Ramírez, S.\ (2018).
  FastAPI framework.
  \url{https://fastapi.tiangolo.com}.

\bibitem{platt}
  Platt, J.\ (1999).
  \emph{Probabilistic Outputs for Support Vector Machines}.
  Advances in Large Margin Classifiers, MIT Press.

\bibitem{qbc}
  Seung, H.\ S., et al.\ (1992).
  \emph{Query by Committee}.
  COLT 1992.

\end{thebibliography}

\end{document}
